\section{Probability}
\subsection{Foundations}\label{sec:foundations}
\subsubsection{Events and probability}
\begin{definition}[Probability measure]\label{def:probability}
We can define the following notation for probability:
\begin{itemize}
    \item We denoted the \textit{state space} or \textit{universal set} as $\Omega$ or $S$.
    \item An \textit{event} is a ``nice" subset of the state space, i.e. an event $A$ fulfils $A\subseteq\Omega$. The set of all events is often denoted by $\mathcal{A},\ \mathcal{E},$ or $\mathcal{F}$.
    \item A \textit{probability measure} $\prob$ is a map taking events as argument with
    \begin{itemize}
        \item $\prob:\ \mathcal{A}\to[0,1]$
        \item $\prob(\Omega)=1$
        \item For a countable index set $I$ with $(A_i)_{i\in I}$ being disjoint events, we have $\prob\brround{\bigsqcup\limits_{i\in I}A_i}=\sum\limits_{i\in I}\prob(A_i)$
    \end{itemize}
\end{itemize}
\end{definition}
Intersection: $A\cap B=\brcurly{\omega\in\Omega:\ \omega\in A\text{ and }\omega\in B}$.
\begin{figure}[htbp]
\centering
    \begin{venndiagram2sets}[labelA={}, labelB={}, labelOnlyA={$A$}, labelOnlyB={$B$}, labelAB={}]
        \fillACapB
    \end{venndiagram2sets}

\caption{Intersection of two events. Existing Venn source; local review only.}\label{fig:intersection}
\end{figure}
\subsubsection{Discrete uniform distribution}
We use the square cups $\sqcup$ to denote the disjoint union of sets. This means that $A_1\sqcup A_2$ implies that the two sets have no overlap and they are disjoint: $A_1\cap A_2=\emptyset$.

We require that $\Omega$ is non empty and finite. If every event $\omega$ is equally likely, we have a discrete uniform distribution. This means that
\[\prob(\{\omega\})=\frac{1}{|\Omega|}\ \forall\omega\in\Omega\]
Then the probability of an event $A$ is
\[\prob(A)=\frac{|A|}{|\Omega|}\]
We call this setting a discrete uniform distribution or uniformly random.


\begin{example}[One die]
For a die, we have $\Omega=\{1,2,3,4,5,6\}$ and $\prob(\{\omega\})=\frac{1}{6}$.
\end{example}
\subsubsection{Inclusion--exclusion}
\begin{theorem}[Two events]\label{thm:inclusion}
Let $A$ and $B$ be two events. Then
\begin{equation}\label{eq:inclusion}
\prob(A\cup B)=\prob(A)+\prob(B)-\prob(A\cap B)
\end{equation}
\end{theorem}
\begin{proof}
Observe that $A\cup B=(A\setminus B)\sqcup(A\cap B)\sqcup(B\setminus A)$ which is a disjoint union. Hence,
\begin{align*}
    \prob(A\cup B)&=\prob(A\setminus B)+\prob(A\cap B)+\prob(B\setminus A)\\
    &=\prob(A\setminus B)+\prob(A\cap B)+\prob(B\setminus A)\\
    &=\prob(A)-\prob(A\cap B)+\prob(A\cap B)+\prob(B)-\prob(A\cap B)\\
    &=\prob((A\setminus B)\sqcup(A\cap B))+\prob((A\cap B)\sqcup(B\setminus A))-\prob(A\cap B)\\
    &=\prob(A)+\prob(B)-\prob(A\cap B)
\end{align*}
\end{proof}
See Definition~\ref{def:probability}, Theorem~\ref{thm:inclusion},
equation~\eqref{eq:inclusion}, Figure~\ref{fig:intersection},
and Section~\ref{sec:foundations}.
